\uuid{mVFe}
\titre{Primitives usuelles}
\auteur{hanani ; mzari}
\organisation{exo7}
\licence{CC BY-NC-SA 4.0}
\datecreate{2020-01-20}
\niveau{L1}
\module{Analyse}
\chapitre{Calcul d'intégrales}
\sousChapitre{Primitives diverses}
\difficulte{3}
\theme{Primitives, primitives usuelles}
\exo7qcmid{lille-2.primitives.042}

\contenu{
\begin{question}
\qclassification{127.11}
\qtags{motcle=primitives usuelles}
Pour $x\in [0,+\infty[$, on pose $t=x^2$. Parmi les assertions suivantes, cocher celles qui sont vraies :
\begin{answers}
\bad{$\displaystyle \int \frac{x\mathrm{d}x}{x^4-x^2-2}=\int \frac{\sqrt{t}\mathrm{d}t}{t^2-t-2}$.}
\good{$\displaystyle \int \frac{x\mathrm{d}x}{x^4-x^2-2}=\frac{1}{2}\int \frac{\mathrm{d}t}{t^2-t-2}$.}
\good{$\forall t\in \Rr\setminus \{-1,2\}$, $\displaystyle \frac{3}{t^2-t-2}=\frac{1}{t-2}-\frac{1}{t+1}$ et $\displaystyle \int \frac{3\, \mathrm{d}t}{t^2-t-2}=\ln\left|\frac{t-2}{t+1}\right|+k$, $k\in \Rr$.}
\bad{$\displaystyle \int \frac{x\mathrm{d}x}{x^4-x^2-2}=\frac{1}{2}\ln\left|\frac{x^2-2}{x^2+1}\right|+k$, $k\in \Rr$.}
\end{answers}
\vskip2mm
\begin{explanations}
Avec $t= x^2$, on a : $\mathrm{d}t=2x\mathrm{d}x$ et
$\displaystyle \int \frac{\mathrm{d}x}{x^4-x^2-2}=\frac{1}{2}\int \frac{\mathrm{d}t}{(t-2)(t+1)}$. Or, $\forall t\in \Rr\setminus \{-1,2\}$, $\displaystyle \frac{3}{t^2-t-2}=\frac{1}{t-2}-\frac{1}{t+1}$. Donc $\displaystyle \int \frac{3\, \mathrm{d}t}{t^2-t-2}=\ln\left|\frac{t-2}{t+1}\right|+k$, $k\in \Rr$, et ensuite :
$\displaystyle \int \frac{x\mathrm{d}x}{x^4-x^2-2}=\frac{1}{6}\ln\left|\frac{x^2-2}{x^2+1}\right|+k$, $k\in \Rr$.
\end{explanations}
\end{question}
}
